\documentclass[a4paper,12pt]{article}




%\abovedisplayshortskip=-12pt plus 3pt
%\setlength{\abovedisplayskip}{-5pt}
%\setlength{\belowdisplayskip}{-5pt}
%\linespread{1.05}
%\voffset=-14cm
%\hoffset=-0.5cm

\usepackage{bm, amsmath, amsfonts, amssymb, url}
\usepackage{titlesec, titletoc}
\usepackage{graphicx,subfigure}
\usepackage{appendix, indentfirst, makeidx, array}
\usepackage{booktabs}
\usepackage{fancyhdr}
\pagestyle{fancy}
\usepackage{fullpage} 

\addtolength{\headsep}{+1cm}
\addtolength{\oddsidemargin}{+0.5cm}
\addtolength{\evensidemargin}{-0.5cm}
%\usepackage[cam,a4,center]{crop}
%\usepackage{geometry}


%\linespread{1.3}

%\fancypagestyle{plain}{
%\fancyhf{}
%\renewcommand{\headrulewidth}{0pt}
%%\fancyfoot[c]{\thepage}
%}


%\addtocontents{toc}{\protect\thispagestyle{empty}}

%%%-------page style ------------

%\setlength{\baselineskip}{20pt plus0.3pt minus0.2pt}
%\setlength{\marginparsep}{0pt}
%\setlength{\marginparwidth}{1.5cm}

%\setlength{\oddsidemargin}{3cm}
%\setlength{\evensidemargin}{3cm}
%%%-------------------------------


\renewcommand{\labelenumi}{(\alph{enumi})}
\renewcommand{\arraystretch}{1.2}
\titleformat{\chapter}[display]{\Large\bf}{\LARGE{Chapter~\thechapter}}{2mm}{}
\titleformat{\section}{\large\bf}{\thesection}{1.5mm}{}
\titleformat{\subsection}{\bfseries\itshape}{\thesubsection}{1.5mm}{}
%self-define symbol
\newcommand{\proposition}[1]{\bigskip\noindent\fbox{\begin{minipage}{0.95\textwidth}~\begin{minipage}{0.95\textwidth}#1\end{minipage}\\\smallskip\end{minipage}}\bigskip\vline\vline\vline}

\newcommand{\excercise}{\addtocounter{section}{+1}
\subsection*{\underline{Ex~\thesection}}}
\newcommand{\stag}[1]{\tag{\arabic{chapter}.\arabic{equation}#1}}

\newcommand{\qdo}[0]{\quad}
\newcommand{\qdt}[0]{\quad\quad}
\newcommand{\qdthr}[0]{\quad\quad\quad}
\newcommand{\qdf}[0]{\quad\quad\quad\quad}
\newcommand{\qdfiv}[0]{\quad\quad\quad\quad\quad}

\newcommand{\norm}[1]{\|#1\|}
%\newcommand{\sdelta}{{\scriptstyle\Delta}}
\newcommand{\sdelta}{\Delta}
\newcommand{\bra}[1]{\langle#1|}
\newcommand{\ket}[1]{|#1\rangle }
\newcommand{\bket}[1]{\Big |#1\Big\rangle }
\newcommand{\braket}[2]{\langle#1|#2\rangle }
\newcommand{\proj}[2]{|#1\rangle\langle#2|}
\newcommand{\expect}[2]{\langle #1\rangle_{#2}}
\newcommand{\Kappa}[0]{\bm{\mathcal{K}}}
\newcommand{\parity}[0]{\bm{\mathcal{P}}}
\newcommand{\parityl}[0]{\bm{p}}
\newcommand{\charge}[0]{\bm{\mathcal{C}}}


\makeindex
%-----------------------------------%

\begin{document}

\abovedisplayskip=-6pt plus 3pt minus 0pt
\belowdisplayskip=8pt plus 3pt minus 0pt

\abovedisplayshortskip=-6pt plus 3pt minus 0pt
\belowdisplayshortskip=8pt plus 3pt minus 0pt

\title{Problem Set 3 (Due: Oct. 12,2011)}
\maketitle

Your homework should be submitted to TA (Chih-Han Lin) at the  beginning of the course. Generally, You have a week to finish the problem set since it was revealed. If you cannot submit it on time, you can still bring your homework to TA's office (L325, Institute of Atomic and Molecular Sciences, in NTU campus) and put it into TA's mailbox which is on the top of the shoe cabinet in L325. Always remember to remind TA to check your homework by e-mail (\url{clin@ltl.iams.sinica.edu.tw}) if you don't submit it in class.

%\abovedisplayskip=0pt %plus 3pt minus 3pt
%\belowdisplayskip=2pt %plus 3pt minus 3pt
%
%\abovedisplayshortskip=0pt %plus 3pt minus 3pt
%\belowdisplayshortskip=2pt %plus 3pt minus 3pt

\excercise
Prove that if $ \mathbf{A} $ and $ \mathbf{B} $ are two operators that both commute with their commutator, then

$$
e^{\mathbf{A}}e^{\mathbf{B}}=e^{\mathbf{A}+\mathbf{B}+[\mathbf{A},\mathbf{B}]/2}.
$$\\[-5pt]

\excercise

If two observables  $ \mathbf{A}_1 $ and $ \mathbf{A}_2 $ are not compatible, but their corresponding operators both commute with the Hamiltonian operator $ \mathbf{H} $, i.e.

$$
[\mathbf{A}_1,\mathbf{H}]=[\mathbf{A}_2,\mathbf{H}]=0.
$$
 
Show that the energy eigenstates are in general degenerate.


\excercise

Consider a one-dimensional Hamiltonian

$$
\mathbf{H}=\frac{1}{2m}\mathbf{P}^2+V(\mathbf{Q})
$$

\noindent and use the fact that the commutator of $ \mathbf{Q} $ and $ [\mathbf{Q},\mathbf{H}] $  is a constant operator, to show that 

$$
\sum_k(E_k-E_s)|Q_{sk}|^2=\frac{\hbar^2}{2m},
$$

\noindent which is referred to as the Thomas-Reiche-Kuhn sum rule,
 where $ Q_{sk}=(\psi_s,\mathbf{Q}\psi_k) $  and  $ \psi_s $ is the eigenstate of $ \mathbf{H} $ with eigenvalue $ E_s $, i.e. $ \mathbf{H}\psi_s=E_s\psi_s $.



\excercise

Prove that if 

 
\begin{align*}
\mathbf{U}(L_3,\theta)=\exp\left[\frac{i}{\hbar}(XP_y-YP_x)\right],
\end{align*}

then

\begin{align*}
\mathbf{U}X\mathbf{U}^{-1}&=X\cos\theta-Y\sin\theta,\\
\mathbf{U}Y\mathbf{U}^{-1}&=X\sin\theta+Y\cos\theta.
\end{align*}

\excercise

Show that if 

\begin{align*}
\mathbf{U}(\vec{L},\vec{\theta})=\exp\left(\frac{i}{\hbar}\vec{\theta}\cdot\vec{L}\right),
\end{align*}

then $\mathbf{U}$ commutes with $\vec{X}\cdot\vec{X}$, $\vec{Y}\cdot\vec{Y}$ and $\vec{Z}\cdot\vec{Z}$. (\textbf{Hint:} Try to show that $\mathbf{U}[X^2+Y^2+Z^2]\mathbf{U}^{-1}=X^2+Y^2+Z^2.$)


\end{document}
